\section{实验结论}

\subsection{贡献}

本实验基于 OpenCV 官方 aloe 双目图像，建立了一个从数据读取、尺寸检查、灰度与同步
降采样、StereoSGBM 匹配，到浮点视差、有效掩码、相对深度和结果归档的完整流程。报告
同时保留了 Mermaid 流程源文件，使处理步骤可以从图示和程序入口两条路径复核。

\subsection{证据}

真实数组统计表明，$355755$ 个像素中有 $261378$ 个满足有效视差条件，占比
$73.4713\%$；有效视差范围为 $16.000$--$105.875$~px，相对深度代理范围为
$0.0094451$--$0.0625$。直接由数组核验 $1/d$ 的最大误差为 $0.0$，说明保存的
相对深度数组与报告公式一致。结果图能够显示叶片、花盆和背景的主要结构，支持
“视差反映相对前后关系”的实验结论。

\subsection{意义与限制}

实验完成了 OpenCV 双目匹配和任意尺度相对深度估计，展示了参数设置、无效区域筛选
以及原始数组与显示图之间的区别。由于 aloe 图像对没有明确配套的焦距 $f$、基线 $B$
和 $\mathbf{Q}$，本实验不能把 \codepath{depth_proxy} 转换为米或厘米，也不输出带物理尺度的
三维点云。遮挡、纹理不足、重复纹理和边界搜索范围会进一步降低有效像素比例。

若要开展绝对深度测量，应先针对同一双目系统完成标定和立体校正，再沿用本实验的
匹配、数组保存和统计核验流程。当前结果的有效范围和限制均来自真实运行数据，未将
未执行的改进步骤写成实验事实。
